\documentclass[../main.tex]{subfiles}

\begin{document}

%% ======================================================================
%%  Executive Summary
%% ======================================================================
\section*{Executive Summary}
\phantomsection
\addcontentsline{toc}{section}{Executive Summary}
\label{sec:executive-summary}

\subsection*{What this report is}

This is a targeted, ecosystem-level technical assessment of safety, security,
and privacy risk in the HDF5 software and data ecosystem, with an evidence
cutoff of August 27, 2026. It is not a certification, attestation, penetration
test, or exhaustive source review, and it expresses no assurance opinion. The
full engagement type, scope, and limitations are stated in
Sections~\ref{sec:engagement-purpose}--\ref{sec:assessment-limitations}.

Ratings apply to specific adverse-event scenarios at specific boundaries, never
to HDF5 as a technology. The rating method is defined once, in
Section~\ref{sec:risk-rubric}, and Appendix~\ref{app:findings-register} is
authoritative for every formal rating in this revision.

\subsection*{Four ecosystem-level judgments}

\begin{enumerate}
    \item \textbf{Historically affected core and official-tool paths repeatedly
          failed to validate file-derived values} --- sizes, counts, offsets,
          message state, layout state, filter and datatype parameters ---
          before using them in trusted C operations. Local fixes can treat a
          known variant; eliminating the class requires invariant-driven
          parsing and continuing regression evidence.
    \item \textbf{Applications can turn documented HDF5 capabilities into
          attack mechanisms} when they process untrusted files while holding
          filesystem, plugin-loading, deserialization, or resource authority
          that the file's originator does not have. The defect is in the
          authority boundary, not necessarily in the parser.
    \item \textbf{Reader acceptance, format conformance, and safe deployment
          are three separate decisions.} Permissive historical acceptance is
          not proof that a file is consistent, and stricter rejection is not
          automatically a regression.
    \item \textbf{Privacy harm can arise during correct, authorized
          operation}, because protecting the payload does not protect
          metadata, structure, derived information, caches, logs, or workflow
          artifacts.
\end{enumerate}

\subsection*{Registered findings at a glance}

Table~\ref{tab:exec-findings-at-a-glance} is a navigation aid. The controlling
records --- scenario statements, scope, evidence, and rating bases --- are in
Appendix~\ref{app:findings-register}; the proposed owners come from
Appendix~\ref{app:phased-roadmap} and are role labels, not present
assignments. Counts here describe bounded scenarios only. They are not
prevalence estimates and must not be aggregated into a score.

\begingroup
\scriptsize
\setlength{\tabcolsep}{2pt}
\renewcommand{\arraystretch}{1.15}
\setlength{\LTleft}{-0.7in}
\setlength{\LTright}{-0.7in}
\begin{longtable}{@{}p{0.145\widetablewidth}>{\raggedright\arraybackslash}p{0.305\widetablewidth}p{0.055\widetablewidth}p{0.095\widetablewidth}p{0.085\widetablewidth}>{\raggedright\arraybackslash}p{0.19\widetablewidth}p{0.065\widetablewidth}@{}}
\caption{Registered findings, ratings, and response routes at a glance}
\label{tab:exec-findings-at-a-glance}\\
\hline
\textbf{ID} & \textbf{Scenario shorthand} & \textbf{Dim.} & \textbf{Tier} &
\textbf{Status} & \textbf{Response route} & \textbf{Owner} \\
\hline
\endfirsthead

\caption[]{Registered findings, ratings, and response routes at a glance,
continued}\\
\hline
\textbf{ID} & \textbf{Scenario shorthand} & \textbf{Dim.} & \textbf{Tier} &
\textbf{Status} & \textbf{Response route} & \textbf{Owner} \\
\hline
\endhead

\hline
\endfoot

\hline
\endlastfoot

\code{APP-SEC-001} &
Affected dataset service followed attacker-selected external raw-storage paths
with worker authority and returned local files. & Sec. & Critical & Rated &
Immediate (24 h from identification) & \code{APP} \\

\code{CORE-SEC-001} &
Malformed on-disk metadata reaches memory-unsafe core-library paths in
historically affected releases. & Sec. & High & Provisional &
Near-term (30 d) & \code{CORE} \\

\code{APP-SEC-002} &
Affected Keras legacy-HDF5 model loading executes embedded code despite
\code{safe_mode=True}. & Sec. & High & Rated & Near-term (30 d) & \code{APP} \\

\code{APP-SEC-003} &
Affected Keras model loading follows attacker-controlled external dataset
references and discloses local files. & Sec. & High & Rated &
Near-term (30 d) & \code{APP} \\

\code{APP-SEC-005} &
TensorFlow 2.17.0 loads HDF5 plugins from an uncontrolled search location,
enabling local privilege escalation. & Sec. & High & Rated &
Near-term (30 d) & \code{APP} \\

\code{CORE-TOOL-001} &
A malicious file triggers the reported \code{h5dump} use-after-free in HDF5
1.14.1-2 and earlier. & Sec. & Mod.--High & Provisional &
Near-term (30 d) & \code{CORE} \\

\code{CORE-PRIV-001} &
Sensitive plaintext metadata is disclosed when a file is shared under
payload-only protection. & Priv. & Mod.--High & Provisional &
Near-term (30 d) & \code{PRIV} \\

\code{CORE-SAFE-001} &
Storage exhaustion during write or close leaves an incomplete or corrupted
file; no library mechanism detects or recovers from it. & Safety & Mod.--High &
Provisional & Near-term (containment) & \code{CORE} \\

\code{CORE-SAFE-002} &
Abrupt process or host termination during write leaves a file corrupted or not
reopenable. & Safety & Mod.--High & Provisional & Near-term (containment) &
\code{CORE} \\

\code{LONG-SAF-001} &
An archived file becomes unreadable because a required filter, VFD, or VOL
connector is gone. Decadal horizon. & Safety & Mod.--High & Provisional &
Planned (recorded departure) & \code{INTEROP} \\

\code{CORE-SEC-002} &
File-controlled work, allocation, traversal, or error paths terminate or
exhaust an affected reader. & Sec. & Moderate & Provisional &
Planned (90 d) & \code{CORE} \\

\code{APP-SEC-004} &
Affected Keras weight loading allocates from an extreme file-declared shape
and exhausts memory. & Sec. & Moderate & Rated & Planned (90 d) & \code{APP} \\

\code{FMT-SAF-001} &
The normative specification is incomplete, ambiguous, or contradicted by the
reference library, so conformance is not determinable. & Safety & Mod.--High &
Provisional & Near-term (adjudicate disputes) & \code{INTEROP} \\

\code{IND-SAF-001} &
HDF5 2.0/2.1 exact-width validation rejects a sufficient but nonminimal chunk
encoding permitted by the specification. & Safety & Moderate & Rated &
Planned (90 d) & \code{INTEROP} \\

\code{IND-SAF-002} &
Older permissive readers accept contradictory PureHDF 2.1.2 chunk metadata
that HDF5 2.0 rejects after upgrade. & Safety & Moderate & Rated &
Planned (90 d) & \code{INTEROP} \\



\hline

\code{CORE-SEC-003} &
Candidate deserializer invariants derived from in-file bytes are enforced only
inside \code{assert(...)} and may be absent from \code{-DNDEBUG} builds. One
instance is confirmed and promoted into \code{CORE-SEC-002}; extent unknown. &
Sec. & \emph{None} & Backlog &
Evidence action (\code{RM-1E}, 30 d target) & \code{CORE} \\

\code{EXT-PRIV-001} &
Extensions can move data and metadata across storage, cache, logging, network,
and executable-plugin boundaries. & Priv. & \emph{None} & Backlog &
Evidence action (\code{RM-2E}) & \code{PRIV} \\

\code{SUP-SAFE-001} &
Downstream builds and workflows can disagree about ABI, linkage, filters,
plugins, VFDs, parallel features, or file locking. & Safety & \emph{None} &
Backlog & Evidence action (\code{RM-2F}) & \code{REL} \\

\end{longtable}
\endgroup

Fifteen formal findings are registered: seven Rated and eight Provisional.
Their tiers comprise one Critical, four High, six Moderate--High ranges, and
four Moderate. Three further observations remain in the evidence-and-scoping
backlog below the rule in the table. \textbf{A backlog record is an evidence
gap, not a determination of Low risk.}

One finding, \code{FMT-SAF-001}, is about the normative artifact rather than
any implementation: where the specification and the reference library disagree
and no procedure exists to rule between them, conformance cannot be
established. Its exposure rises with the number of independent
implementations, because each unresolved ambiguity is resolved independently by
every one of them.

Three further findings --- \code{CORE-SAFE-001}, \code{CORE-SAFE-002}, and
\code{LONG-SAF-001} --- are write-side and lifecycle safety risks that an
earlier revision of this report either left unrated or did not register at all.
They are called out here because they are the conditions HDF5 users encounter
most often in ordinary operation, and because the rest of this report is
organised around a different question: hostile input arriving from outside,
rather than data failing to survive.

\subsection*{Where the priority actually falls}

Response urgency comes from Table~\ref{tab:findings-response-routing} and
nowhere else. Chapter order, recommendation numbering, CVE counts, publicity,
effort, and dependencies do not create or revise it.

\begin{itemize}
    \item \textbf{The single Immediate route belongs to downstream service
          operators, not to the reference-library maintainers.} It is
          conditional: it applies to any service that still automatically
          follows file-selected external resources while holding access to
          sensitive local files. The affected July 2026 deployment is
          reported remediated, and this report assigns no current residual
          tier to it.
    \item \textbf{For The HDF Group specifically, there is no Immediate
          action in this revision.} The near-term core obligations are the
          deployment census (\code{RM-1A}), the affected-version treatment of
          \code{CORE-SEC-001} and \code{CORE-TOOL-001} (\code{RM-1B}), and the
          30-day assert-site catalog (\code{RM-1E}).
    \item \textbf{Two response clocks are distinct.} A finding's clock starts
          when an affected deployment or equivalent exposure is
          \emph{identified}. The census and assert-catalog clocks start when
          the roadmap is \emph{adopted}. Adoption must set due dates for both
          inventories, so that a response clock cannot stay latent through
          unbounded non-discovery.
\end{itemize}

\subsection*{If only a few things can be funded}

The roadmap in Appendix~\ref{app:phased-roadmap} is scoped to risk, not to
budget. Section~\ref{sec:roadmap-minimum-viable} identifies a
capacity-constrained subset that preserves the response clocks and delivers
the highest control coverage per unit of effort. It is a funding aid and
explicitly not a risk rank.

\subsection*{What this report does not establish}

No proposed mitigation is shown to be adopted, implemented, or effective.
Every roadmap package is \emph{Proposed / implementation not verified} at
publication, and no residual-risk acceptance is recorded. Current installed
populations, exact fix coverage across supported releases, and
deployment-specific control effectiveness remain incomplete. Unless a record
expressly states otherwise, findings rest on document, artifact, and
architecture examination rather than independent reproduction. The absence of
a reported finding does not establish the absence of defects.

\end{document}
